\section{Discussion}
Across Denmark, the UK, and the US, judgments of EDC-related exposure scenarios revealed common patterns across psychometric dimensions and perceived exposure pathways, alongside more variable associations with individual predictors. Participants assessed these scenarios within the affective-dread and cognitive-deliberate dimensions of the psychometric paradigm \citep{fischhoff1978safe,slovic1987perception}. They relied on a stable exposure pathway hierarchy, perceiving consumption-related exposures as riskier than inhalation and dermal application. Associations with individual predictors were less uniform, although self-reported EDC knowledge showed consistent associations across countries and specific risk outcomes. Beyond describing EDC-related risk perceptions, these findings suggest three broader contributions. First, psychometric risk structures seem to extend to chronic and diffuse chemical hazards embedded in everyday life. Second, exposure pathways may serve as a psychologically meaningful organizing cue in EDC judgments. Third, cognitive variables may provide a more robust basis for understanding EDC risk perceptions than broad demographic profiles.

\subsection{The Affective-Cognitive Structure of EDC Risk Perception}
EDC-related risk perceptions reproduced the canonical two-dimensional structure described by the psychometric paradigm \citep{fischhoff1978safe, slovic1987perception}. Across Denmark, the UK, and the US, exposure scenarios clustered along an affective-dread dimension and a cognitive-deliberate dimension, indicating that participants evaluated these risks in a structured manner. Thus, our findings are consistent with prior psychometric research showing that lay risk judgments are organized by a limited number of broad dimensions reflecting both affective and deliberative evaluations \citep{slovic1987perception, slovic2012perception}. They also align with dual-process accounts suggesting that risk judgments often reflect an interaction between rapid, affect-laden impressions and more reflective reasoning \citep{kahneman2011thinking, finucane2000affect}.

Our findings suggest that the psychometric paradigm remains informative even for hazards that are chronic, cumulative, and embedded in ordinary consumption and household practices. Much of the classic psychometric literature focused on dramatic technological or environmental hazards \citep{fischhoff1978safe, slovic1987perception}, whereas EDC exposures are more diffuse, probabilistic, and difficult to link to direct personal experience. The emergence of the same broad structure across all three countries suggests that psychometric dimensions may generalize beyond the acute and catastrophic hazards that often dominated early psychometric risk research. This interpretation also aligns with recent work showing that lay understandings of chemicals often diverge from formal toxicological distinctions and are shaped by familiarity, controllability, and partial understandings of hazard and exposure \citep{bearth2019lay, saleh2019chemophobia}.

Although the same two dimensions emerged in all three countries, their relative salience differed. The cognitive-deliberate dimension accounted for more variance in Denmark, whereas the affective-dread dimension was more prominent in the UK and US. This finding suggests that cross-national variation may operate less through fundamentally different perceptual structures than through differences in the relative weighting of shared dimensions. This interpretation is consistent with more recent work emphasizing that risk judgments are shaped by communication environments, institutional credibility, and trust rather than emerging in a social vacuum \citep{balog2020evolving, kasperson2022social, bearth2022social}.

The psychometric profiles also largely aligned with our preregistered expectations while revealing important boundary conditions. Across countries, EDC-related exposure scenarios were generally perceived as delayed rather than immediate risks and as substantially better understood by scientists than by the public, supporting the characterization of EDCs as “slow-burn” hazards accompanied by a perceived public–expert knowledge asymmetry. Pesticide-related scenarios were also viewed as relatively less novel than other EDC-related exposures. Judgments of chronicity, dread, and fatality clustered near scale midpoints, suggesting that EDC-related risks were not construed as catastrophic, highly dreaded, or certainly fatal.

Perceptions of voluntariness and controllability followed meaningful patterns across exposure pathways. Inhalation-related exposures were generally viewed as more involuntary and less controllable, whereas dermal application exposures were viewed as more voluntary and controllable. Consumption-related exposures occupied an intermediate position. These patterns suggest that EDCs are not represented as a single psychological category; rather, perceptions vary depending on the exposure pathway and the degree of perceived personal agency involved. Theoretically, this finding suggests that psychometric attributes such as controllability and voluntariness may depend partly on exposure context rather than reflecting stable properties of chemical hazards themselves.

\subsection{Exposure Pathways as an Organizing Cues in EDC Risk Perception}
A stable exposure pathway hierarchy emerged across all three countries. This pattern provided partial support for our preregistered hypothesis, as consumption-related exposures were consistently perceived as riskier than inhalation and dermal application exposures. Contrary to our expectations, inhalation-related scenarios were judged riskier than dermal application scenarios. This finding diverges from research on contagion, which suggests that direct bodily contact often amplifies perceived contamination and risk \citep{rozin1986operation, rozin1992magical}. At the same time, it aligns with emerging EDC-specific research suggesting that lay understandings of chemical risks are shaped by perceived exposure sources and pathways and perceived control over exposure \citep{kelly2020public, boronow2025people}. Research on public perceptions of chemical exposure further shows that people are generally aware of ingestion, inhalation, and dermal exposure pathways while expressing lower concern about dermal absorption \citep{matisane2022citizens}. Together, these findings suggest that judgments of chemical exposure may depend less on physical contact and more on how visible, controllable, and mentally accessible an exposure pathway appears. Theoretically, these findings suggest that exposure pathways function as a stable organizing cue in lay judgments of chronic chemical risk across national contexts and may complement classic psychometric dimensions when risks are diffuse, cumulative, and embedded in everyday life.

\subsection{Individual Predictors of EDC Risk Perception}
Across countries and outcomes, self-reported EDC knowledge was the most stable predictor particularly of perceived severity, whereas many sociodemographic effects were weaker, more variable, or outcome-specific. This pattern is consistent with recent synthesis work identifying knowledge as one of the most robust cognitive correlates of perceived EDC risk \citep{pravednikov2024main}. It also suggests that greater perceived knowledge may heighten awareness of the diffuse, long-term, and often difficult-to-avoid nature of EDC exposures rather than reassure individuals about them.

Gender, age, parenthood, education, and income showed significant associations, but these patterns were not stable enough to support a single universal profile of the "high-risk perceiver". Theoretically, this suggests that EDC risk perception cannot be captured by demographic characteristics alone. Perceptions depend on interactions among cognitive variables, institutional context, and the specific dimension of perceived risk being studied. Thus, our findings converge on and refine the broader risk-perception literature, which has long documented demographic variation in perceived risk \citep{flynn1994gender, finucane2000affect}, while increasingly recognizing the importance of context, framing, and institutional trust \citep{siegrist2021trust, balog2020evolving}.

Prior familiarity with EDCs reduced perceived severity in Denmark, suggesting that greater familiarity may normalize rather than intensify concerns in some contexts. Trust also operated differently across countries, reducing perceived overall health risk to EDCs in Denmark while increasing perceived severity in the US and UK. These findings suggest that trust does not simply attenuate perceived risk. Depending on the institutional setting, it may alternatively function as reassurance or signal that a risk deserves serious attention. This interpretation aligns with recent work on the Social Amplification of Risk Framework and risk communication, which emphasizes that trust is embedded within broader communication environments rather than functioning as a simple factor that uniformly increases or decreases perceived risk \citep{bearth2022social, kasperson2022social}. 

\subsection{Distinguishing Perceived Susceptibility, Severity, and Overall Health Risk}
Our findings suggest that perceived susceptibility, severity, and overall health risk should not be treated as interchangeable indicators of EDC risk perception. Although perceived susceptibility and perceived severity were moderately associated, participants consistently rated susceptibility higher than severity. Moreover, pathway-specific rankings of perceived health risk had weak associations with perceived overall health risk judgments. Thus, people may perceive EDC-related illness as relatively likely without necessarily judging its consequences as equally severe. Conversely, stronger judgments of health risks associated with particular exposure pathways did not straightforwardly translate into higher overall perceptions of EDC-related health risk.

This pattern is broadly consistent with health-belief approaches that distinguish perceived susceptibility from perceived severity as related but conceptually distinct components of perceived threat \citep{rosenstock1974historical, janz1984health}. At the same time, our findings suggest that overall judgments about EDC-related health risk may capture something partially different from either construct. Rather than reflecting a simple aggregation of pathway-specific beliefs, these broader judgments may incorporate impressions of ubiquity, uncertainty, uncontrollability, and difficulty of avoidance. This interpretation aligns with recent work suggesting that lay understandings of chemicals often combine hazard, ubiquity, uncertainty, and controllability into a more general sense that chemicals are "everywhere" or difficult to escape \citep{bearth2019lay, saleh2019chemophobia, boronow2025people}. Theoretically, these findings suggest that different dimensions of EDC risk perception should be distinguished more clearly in both theory and measurement.

\subsection{Policy Recommendations}

Policy efforts to reduce EDC exposure can combine regulatory measures with support for exposure-reducing behavior, although opportunities for individual action differ across exposure contexts. Based on our findings, we draw five recommendations.

First, communication could be organized around exposure pathways rather than individual chemicals because risk judgments clustered primarily by pathway (Figure 3). Existing consumer-facing tools provide examples of actionable communication in the dermal domain, including cosmetic-screening applications and personal-care guides \citep{forbrugerraadet_kemiluppen,bcpp_personalcare}, illustrate this approach. Yet, dermal exposures were judged both most controllable and least risky, suggesting that communication in this domain may need to make the potential health relevance of dermal exposures more salient, while similarly actionable guidance could be developed for pathways that were perceived as riskier.

Second, where exposure is perceived as involuntary and uncontrollable, particularly inhalation, communication should emphasize feasible actions that can increase agency, such as ventilation and HEPA filter vacuuming \citep{feldscher2026chemicals}. Where individual control is limited, upstream measures may be more appropriate, including regulation of exposure sources such as flame retardants in furniture \citep{volz2025proposal}.

Third, public agencies should translate established evidence into accessible, actionable guidance. Participants consistently perceived EDC risks as better understood by scientists than by the public, while greater self-reported knowledge was associated with higher perceived severity across all three countries. More information should, therefore, not be assumed to reassure audiences; pairing evidence with concrete actions may be more useful.

Fourth, outreach may be better tailored to knowledge and institutional trust than to broad demographic profiles. No stable sociodemographic "high-risk perceiver" emerged across countries and outcomes, whereas knowledge showed the most consistent associations and trust showed country-specific associations. Agencies could, therefore, adapt the depth and framing of information to audience knowledge and trust conditions.

Finally, multinational and EU-level communication strategies should allow for adaptation across national contexts. Cognitive-deliberate considerations were more prominent in Denmark, whereas affective-dread considerations were more prominent in the UK and US. Communication may, therefore, benefit from emphasizing knowledge and controllability where deliberative judgments are more salient and addressing affective concerns where dread is more prominent.

\subsection{Limitations}
First, because the study is cross-sectional and relies on self-reported measures, causal interpretations are limited. For example, we cannot determine whether greater knowledge increases perceived overall health risk and severity, whether heightened concern motivates information seeking (thereby increasing knowledge), or whether both are driven by unmeasured third variables such as media attention or political identity.

Second, structured context effects may have been introduced by the study design. Participants were first provided with a definition of EDCs and were required to correctly identify five EDCs before proceeding with the survey. While this procedure strengthens internal validity by ensuring a minimal and shared understanding of the concept, it also implies that the reported risk perceptions reflect judgments formed after an informational primer rather than baseline, naturally occurring public beliefs.

A third limitation is that participants' responses should not be interpreted as direct measures of stable, pre-existing attitudes toward every EDC-related exposure scenario. As \citet{slovic1995construction} argues, preferences may be constructed during elicitation rather than simply revealed from stable prior evaluations. This concern may be particularly relevant in the present study because participants first received a definition of EDCs and completed an identification task before evaluating exposure scenarios. While this procedure strengthened a shared minimum understanding of EDCs across participants, it also means that reported judgments may partly reflect responses formed within the study context rather than naturally occurring public beliefs. Some evaluations, particularly for less familiar chemical labels or exposure contexts, may therefore have been shaped during participation itself. However, judgments formed during elicitation may still reflect coherent and theoretically meaningful patterns. In the present study, responses converged with theoretically expected findings and prior EDC risk research, suggesting coherent structure in the elicited responses. Accordingly, the findings should be interpreted as describing how EDC-related risks are represented under a standardized information and response context rather than as evidence of stable pre-existing attitudes toward each exposure scenario. 

Fourth, the present study does not link risk perceptions to downstream behaviors or to objective exposure indicators. This limitation is important because prior research shows that perceived contamination, uncertainty, and institutional responses can shape well-being, stress, and community reactions in ways that are not fully captured by psychometric dimensions alone \citep{wakefield2000environmental, sullivan2021chronic}. 

Fifth, the present study focused on perceived risks and did not examine perceived benefits or trade-offs associated with EDC-related applications. This is important as some EDCs serve important technical and functional roles and may be difficult to replace \citep{cousins2019concept,jacobs2015alternatives}.

Finally, chemical labels varied in specificity: some were technical names or acronyms (e.g., BPA, phthalates), whereas others were broader categories (e.g., pesticides). Although selected to balance scientific precision with cross-national recognizability, this variation may have affected perceived naturalness, syntheticity, familiarity, or processing fluency and thereby influenced risk judgments beyond differences in exposure pathway. Because these characteristics were neither manipulated nor measured, their effects cannot be isolated from chemical identity, product context, or prior reputation. Our exploratory analysis (S9) found that technical labels were not uniformly perceived as riskier than pesticides, suggesting a simple processing-fluency explanation is unlikely, but it does not rule out naming effects.

\subsection{Future Research}
These limitations, together with our findings, raise four directions for future research. Most directly, future work should examine what accounts for the observed cross-national differences, including whether regulatory culture, institutional trust, and communication environments account for the greater salience of cognitive-deliberate considerations in Denmark and affective-dread considerations in the UK and US \citep{balog2020evolving, bearth2022social, clahsen2019countries}.

Second, what produces the exposure-pathway hierarchy? Contagion accounts emphasize contact with potential contaminants \citep{rozin1986operation, rozin1992magical, rozin2023magical}, while pathogen-avoidance research suggests that perceived bodily entry also matters \citep{white2025disgust}. Whether these intuitions extend to chemical exposure remains unclear, but they may help explain why participants consistently perceived dermal exposures as less risky, consistent with previous survey evidence \citep{matisane2022citizens}. Perceived control alone, however, cannot explain the hierarchy, since consumption was judged riskiest despite being perceived as more controllable than inhalation. Future research should examine whether perceived bodily entry, beliefs about the skin as a barrier, familiarity, and perceived control jointly account for differences in risk perception across exposure pathways.

Third, because judgments clustered primarily by exposure pathway, future research could test whether pathway-based communication improves understanding and protective behavior, and at what level of aggregation such grouping remains meaningful. This could be compared with established approaches such as dose-response framing \citep{bearth2021chemophobia} and tested across pathways differing in perceived controllability.

Fourth, future work should determine which risk judgments predict behavior. Participants rated susceptibility higher than severity, and pathway-specific rankings were only weakly related to overall risk judgments, consistent with evidence that responses to chemical presence may differ from judgments of health consequences \citep{jansen2020all}. Perceived benefits may also compete with risk, particularly for dermal products such as sunscreen. Longitudinal and behavioral studies linking these judgments to avoidance behavior and biomarker-assessed exposure could establish which perceptions actually predict exposure-reducing behavior.